% ═══════════════════════════════════════════════════════════════
% LERNBLATT: Describir fotos
% Klasse 8 · Spanisch
% ═══════════════════════════════════════════════════════════════

\documentclass[10pt, a4paper]{article}

\usepackage[utf8]{inputenc}
\usepackage[T1]{fontenc}
\usepackage[ngerman]{babel}

\usepackage{helvet}
\renewcommand{\familydefault}{\sfdefault}
\usepackage{palatino}
\newcommand{\seriftext}[1]{{\fontfamily{ppl}\selectfont #1}}

\usepackage{microtype}
\usepackage[margin=1.5cm, top=1.8cm, bottom=1.3cm]{geometry}
\usepackage{enumitem}
\usepackage{tabularx}
\usepackage{booktabs}
\usepackage{array}
\usepackage{multicol}
\usepackage{tikz}

\usepackage{fancyhdr}
\usepackage{lastpage}

% ═══════════════════════════════════════════════════════════════
% VARIABLEN
% ═══════════════════════════════════════════════════════════════

\newcommand{\thema}{Describir fotos}
\newcommand{\untertitel}{Bildbeschreibung auf Spanisch}
\newcommand{\klasse}{8}
\newcommand{\lernziele}{Presente continuo · Bildaufbau · Positionen · Personen beschreiben}

% ═══════════════════════════════════════════════════════════════
% FARBEN
% ═══════════════════════════════════════════════════════════════

\usepackage{xcolor}

\definecolor{colorrojo}{HTML}{C62828}
\definecolor{colorazul}{HTML}{1565C0}
\definecolor{colorverde}{HTML}{2E7D32}
\definecolor{colorgris}{HTML}{757575}
\definecolor{colornegro}{HTML}{222222}
\definecolor{hellgrau}{HTML}{F5F5F5}
\definecolor{rahmen}{HTML}{CCCCCC}
\definecolor{akzent}{HTML}{666666}

% ═══════════════════════════════════════════════════════════════
% BOXEN
% ═══════════════════════════════════════════════════════════════

\usepackage{tcolorbox}
\tcbuselibrary{skins}

\newtcolorbox{lernzielbox}{
    colback=hellgrau, colframe=hellgrau, boxrule=0pt, arc=0pt,
    left=8pt, right=8pt, top=6pt, bottom=6pt,
    before skip=4pt, after skip=6pt
}

\newtcolorbox{grammatikbox}[1][]{
    colback=white, colframe=white, boxrule=0pt, arc=0pt,
    left=8pt, right=6pt, top=5pt, bottom=5pt,
    borderline west={3pt}{0pt}{akzent},
    before skip=3pt, after skip=3pt
}

\newtcolorbox{merkbox}[1][]{
    colback=hellgrau, colframe=hellgrau, boxrule=0pt, arc=0pt,
    left=8pt, right=8pt, top=6pt, bottom=6pt,
    borderline north={2pt}{0pt}{colornegro},
    before skip=4pt, after skip=4pt
}

\newtcolorbox{vokabelbox}[1][]{
    colback=white, colframe=rahmen, boxrule=0.5pt, arc=0pt,
    left=6pt, right=6pt, top=5pt, bottom=4pt,
    before skip=3pt, after skip=3pt
}

\newtcolorbox{tippbox}{
    colback=white, colframe=white, boxrule=0pt, arc=0pt,
    left=6pt, right=4pt, top=2pt, bottom=2pt,
    borderline west={2pt}{0pt}{akzent}
}

% ═══════════════════════════════════════════════════════════════
% BEFEHLE
% ═══════════════════════════════════════════════════════════════

\newcommand{\es}[1]{\seriftext{\textbf{#1}}}
\newcommand{\de}[1]{\textit{\textcolor{akzent}{#1}}}
\newcommand{\gram}[1]{\textsc{#1}}
\newcommand{\abmark}[1]{{\footnotesize\textcolor{akzent}{$\rightarrow$ AB #1}}}
\newcommand{\abschnitt}[2]{\noindent\textbf{\large #1. #2}}
\newcommand{\stamm}[1]{\textcolor{colorrojo}{#1}}

\setlength{\parindent}{0pt}
\setlength{\parskip}{4pt}
\setlength{\columnsep}{18pt}

% ═══════════════════════════════════════════════════════════════
% KOPF- UND FUSSZEILE
% ═══════════════════════════════════════════════════════════════

\pagestyle{fancy}
\fancyhf{}
\renewcommand{\headrulewidth}{0.5pt}
\renewcommand{\footrulewidth}{0pt}
\lhead{\footnotesize\textsc{Spanisch} · Klasse \klasse}
\rhead{\footnotesize\thema}
\cfoot{\footnotesize\thepage\,/\,\pageref{LastPage}}

% ═══════════════════════════════════════════════════════════════
% DOKUMENT
% ═══════════════════════════════════════════════════════════════

\begin{document}

% --- TITEL ---
\begin{center}
    {\LARGE\bfseries \thema}\\[2pt]
    {\small \untertitel}
\end{center}

\vspace{2pt}

Stell dir vor, du zeigst einem spanischen Freund ein Foto aus deinem Urlaub.
Was siehst du? Wer macht was? Heute lernst du, Bilder auf Spanisch zu beschreiben!

\vspace{2pt}

\begin{lernzielbox}
\textbf{Das lernst du:}\quad \lernziele
\end{lernzielbox}

% ═══════════════════════════════════════════════════════════════
% ZWEI SPALTEN
% ═══════════════════════════════════════════════════════════════

\begin{multicols}{2}

% ---------------------------------------------------------------
% G1: PRESENTE CONTINUO
% ---------------------------------------------------------------
\abschnitt{G1}{El presente continuo} \de{-- Verlaufsform}

\vspace{3pt}

\begin{grammatikbox}
\textbf{Bildung:} \es{estar} + \es{gerundio}

\vspace{3pt}

\de{Wie im Englischen:}\\
\textit{I \textbf{am reading}} = \es{Estoy leyendo}

\vspace{4pt}

{\small
\begin{tabular}{@{}ll@{}}
\textbf{estar:} & estoy, estás, está, estamos, estáis, están \\
\end{tabular}}

\vspace{4pt}

\textbf{Gerundio-Endungen:}\\
{\small
\begin{tabular}{@{}lll@{}}
\es{-ar} & $\rightarrow$ & \es{-ando} \de{(hablar → hablando)} \\
\es{-er/-ir} & $\rightarrow$ & \es{-iendo} \de{(comer → comiendo)} \\
\end{tabular}}
\end{grammatikbox}

\vspace{2pt}

\begin{tippbox}
\footnotesize\textit{Achtung:} \es{leer → leyendo}, \es{oír → oyendo} \de{(e/o zwischen Vokalen wird zu y)}
\end{tippbox}

\abmark{G1, A1}

\vspace{6pt}

% ---------------------------------------------------------------
% V1: BILDAUFBAU
% ---------------------------------------------------------------
\abschnitt{V1}{Las partes de la imagen} \de{-- Bildaufbau}

\vspace{3pt}

\begin{vokabelbox}
{\small
\begin{tabular}{@{}ll@{}}
\es{en primer plano} & \de{im Vordergrund} \\
\es{en segundo plano} & \de{im Mittelgrund} \\
\es{al fondo} & \de{im Hintergrund} \\
\es{a la izquierda} & \de{links} \\
\es{a la derecha} & \de{rechts} \\
\es{en el centro} & \de{in der Mitte} \\
\es{arriba} & \de{oben} \\
\es{abajo} & \de{unten} \\
\end{tabular}}
\end{vokabelbox}

\abmark{A2}

\vspace{6pt}

% ---------------------------------------------------------------
% V2: POSITIONEN
% ---------------------------------------------------------------
\abschnitt{V2}{Las posiciones} \de{-- Positionen}

\vspace{3pt}

\begin{vokabelbox}
{\small
\begin{tabular}{@{}ll@{}}
\es{estar de pie} & \de{stehen} \\
\es{estar sentado/a} & \de{sitzen} \\
\es{estar tumbado/a} & \de{liegen} \\
\es{estar apoyado/a} & \de{sich anlehnen} \\
\es{estar agachado/a} & \de{hocken, geduckt sein} \\
\end{tabular}}
\end{vokabelbox}

\begin{tippbox}
\footnotesize\textit{Merke:} Endung passt zum Geschlecht:\\
\es{El hombre está sentado.} / \es{La mujer está sentada.}
\end{tippbox}

\abmark{A3}

\columnbreak

% ---------------------------------------------------------------
% V3: PERSONENBESCHREIBUNG (REPASO)
% ---------------------------------------------------------------
\abschnitt{V3}{Describir personas} \de{-- Repaso}

\vspace{3pt}

\begin{vokabelbox}
{\small
\begin{tabular}{@{}ll@{}}
\multicolumn{2}{@{}l@{}}{\textbf{Größe:}} \\
\es{alto/a} / \es{bajo/a} & \de{groß / klein} \\[3pt]
\multicolumn{2}{@{}l@{}}{\textbf{Haare:}} \\
\es{el pelo largo/corto} & \de{langes/kurzes Haar} \\
\es{el pelo rubio/moreno} & \de{blondes/braunes Haar} \\
\es{el pelo negro/pelirrojo} & \de{schwarzes/rotes Haar} \\
\es{el pelo liso/rizado} & \de{glattes/lockiges Haar} \\[3pt]
\multicolumn{2}{@{}l@{}}{\textbf{Augen:}} \\
\es{los ojos azules/verdes} & \de{blaue/grüne Augen} \\
\es{los ojos marrones} & \de{braune Augen} \\[3pt]
\multicolumn{2}{@{}l@{}}{\textbf{Sonstiges:}} \\
\es{llevar gafas} & \de{eine Brille tragen} \\
\es{joven / mayor} & \de{jung / älter} \\
\end{tabular}}
\end{vokabelbox}

\abmark{A5}

\vspace{6pt}

% ---------------------------------------------------------------
% V4: ORTE
% ---------------------------------------------------------------
\abschnitt{V4}{Los lugares} \de{-- Orte}

\vspace{3pt}

\begin{vokabelbox}
{\small
\begin{tabular}{@{}ll@{}}
\multicolumn{2}{@{}l@{}}{\textbf{Im Park:}} \\
\es{el parque} & \de{der Park} \\
\es{el banco} & \de{die Parkbank} \\
\es{el árbol} & \de{der Baum} \\
\es{el césped} & \de{der Rasen} \\
\es{la fuente} & \de{der Brunnen} \\
\es{el camino} & \de{der Weg} \\[3pt]
\multicolumn{2}{@{}l@{}}{\textbf{In der Stadt:}} \\
\es{la calle} & \de{die Straße} \\
\es{la plaza} & \de{der Platz} \\
\es{el edificio} & \de{das Gebäude} \\
\es{la tienda} & \de{das Geschäft} \\[3pt]
\multicolumn{2}{@{}l@{}}{\textbf{Im Café:}} \\
\es{la cafetería} & \de{das Café} \\
\es{la terraza} & \de{die Terrasse} \\
\es{la mesa} & \de{der Tisch} \\
\es{la silla} & \de{der Stuhl} \\
\end{tabular}}
\end{vokabelbox}

\abmark{A5}

\end{multicols}

% ═══════════════════════════════════════════════════════════════
% SEITE 2
% ═══════════════════════════════════════════════════════════════
\newpage

% ---------------------------------------------------------------
% V5: VERBEN (GERUNDIO)
% ---------------------------------------------------------------
\abschnitt{V5}{Verbos para describir} \de{-- Verben zum Beschreiben}

\vspace{3pt}

\begin{merkbox}
{\small
\begin{multicols}{3}
\begin{tabular}{@{}ll@{}}
\es{caminar} & \de{laufen} \\
\es{pasear} & \de{spazieren} \\
\es{correr} & \de{rennen} \\
\es{andar} & \de{gehen} \\
\end{tabular}

\columnbreak

\begin{tabular}{@{}ll@{}}
\es{mirar} & \de{schauen} \\
\es{hablar} & \de{sprechen} \\
\es{reír} & \de{lachen} \\
\es{sonreír} & \de{lächeln} \\
\end{tabular}

\columnbreak

\begin{tabular}{@{}ll@{}}
\es{comer} & \de{essen} \\
\es{beber} & \de{trinken} \\
\es{leer} & \de{lesen} \\
\es{jugar} & \de{spielen} \\
\end{tabular}
\end{multicols}}
\end{merkbox}

\abmark{A1}

\vspace{6pt}

% ---------------------------------------------------------------
% R1: REDEMITTEL
% ---------------------------------------------------------------
\abschnitt{R1}{Redemittel} \de{-- Satzanfänge für Bildbeschreibungen}

\vspace{3pt}

\begin{multicols}{2}

\begin{grammatikbox}
\textbf{Einstieg:}\\
\es{En la foto veo...} \de{Auf dem Foto sehe ich...}\\
\es{La imagen muestra...} \de{Das Bild zeigt...}\\
\es{Se trata de...} \de{Es handelt sich um...}

\vspace{4pt}

\textbf{Bildaufbau:}\\
\es{En primer plano hay...} \de{Im Vordergrund gibt es...}\\
\es{Al fondo se ve...} \de{Im Hintergrund sieht man...}\\
\es{A la izquierda/derecha...} \de{Links/Rechts...}
\end{grammatikbox}

\columnbreak

\begin{grammatikbox}
\textbf{Aktionen beschreiben:}\\
\es{Una persona está + gerundio...}\\
\de{Eine Person ist dabei zu...}\\
\es{Hay alguien que está...}\\
\de{Es gibt jemanden, der gerade...}

\vspace{4pt}

\textbf{Vermutungen:}\\
\es{Parece que...} \de{Es scheint, dass...}\\
\es{Probablemente...} \de{Wahrscheinlich...}\\
\es{Creo que...} \de{Ich glaube, dass...}
\end{grammatikbox}

\end{multicols}

\abmark{A4, A5}

\vspace{6pt}

% ---------------------------------------------------------------
% MUSTERBESCHREIBUNG
% ---------------------------------------------------------------
\abschnitt{}{Ejemplo: Descripción de una foto}

\vspace{3pt}

\begin{merkbox}
\es{En la foto veo un parque en una ciudad. En primer plano hay un banco donde una mujer está sentada. Ella tiene el pelo largo y rubio. Está leyendo un libro. A la izquierda, un hombre joven está paseando con su perro. Él lleva gafas y tiene el pelo corto y moreno. Al fondo se ven varios edificios altos. En el centro de la imagen hay una fuente bonita. A la derecha, dos niños están jugando en el césped. Parece que es un día soleado de verano.}
\end{merkbox}

\abmark{A4}

\vspace{6pt}

% ---------------------------------------------------------------
% TIPPS
% ---------------------------------------------------------------
\begin{grammatikbox}
\textbf{5 Tipps für eine gute Bildbeschreibung:}

\vspace{3pt}

{\small
\begin{enumerate}[leftmargin=1.5em, itemsep=1pt]
    \item \textbf{Struktur:} Beginne mit dem Überblick, dann Details (Vordergrund → Hintergrund)
    \item \textbf{Presente continuo:} Beschreibe Aktionen mit \es{estar + gerundio}
    \item \textbf{Positionen:} Nutze \es{estar de pie, sentado, tumbado...}
    \item \textbf{Personen:} Beschreibe Aussehen (Haare, Augen, Größe)
    \item \textbf{Vermutungen:} Verwende \es{parece que...} für Interpretationen
\end{enumerate}}
\end{grammatikbox}

\end{document}
